\documentclass[10pt,a4paper]{article}
\usepackage{fontspec}
\usepackage[english,russian]{babel}
\babelfont{rm}{PT Serif}
\babelfont{sf}{PT Sans}
\babelfont{tt}[Scale=0.88]{PT Mono}
\usepackage[a4paper,top=20mm,bottom=22mm,left=21mm,right=21mm]{geometry}
\usepackage{amsmath}
\usepackage{booktabs}
\usepackage{array}
\usepackage{enumitem}
\usepackage{titlesec}
\usepackage{caption}
\usepackage{tikz}
\usetikzlibrary{arrows.meta}
\usepackage{xcolor}
\usepackage[hidelinks]{hyperref}
\usepackage{microtype}
\usepackage{multicol}

\setlength{\parindent}{0pt}
\setlength{\parskip}{3pt plus 1pt}
\setlength{\emergencystretch}{4em}

\titleformat{\section}{\normalfont\fontsize{12.5}{15}\bfseries}{\thesection.}{0.5em}{}
\titlespacing*{\section}{0pt}{11pt plus 2pt minus 2pt}{4pt}
\titleformat{\subsection}{\normalfont\fontsize{11}{13}\bfseries}{\thesubsection.}{0.5em}{}
\titlespacing*{\subsection}{0pt}{7pt plus 2pt minus 1pt}{2pt}

\setlist[enumerate]{leftmargin=2.0em,itemsep=1pt,topsep=3pt,parsep=0pt}
\setlist[itemize]{leftmargin=1.6em,itemsep=1pt,topsep=3pt,parsep=0pt}
\captionsetup{font=small,labelfont=bf,labelsep=period,skip=4pt}
\newcolumntype{L}[1]{>{\raggedright\arraybackslash}p{#1}}
\setlength{\textfloatsep}{9pt plus 2pt minus 2pt}
\setlength{\floatsep}{9pt plus 2pt minus 2pt}

\pagestyle{plain}

\tikzset{
  box/.style={draw,rectangle,thick,align=center,inner sep=3pt,font=\footnotesize,minimum width=26mm,minimum height=8mm},
  zone/.style={draw,dashed,rounded corners=2pt},
  zl/.style={font=\scriptsize\itshape,anchor=north west,inner sep=2pt},
  lbl/.style={font=\scriptsize,align=center,fill=white,inner sep=1pt},
  arr/.style={-{Latex[length=1.8mm]},thick},
}

\begin{document}

\begin{center}
{\fontsize{17}{21}\selectfont\bfseries Swarma: платформа\\[2pt] для роёв автономных агентов\par}
\vspace{6pt}
{\normalsize Рантайм на машинах владельца, подключаемые рабочие процессы,\\ протокол независимых владельцев и единая личность участника\par}
\vspace{10pt}
{\normalsize Алексей Григорьев, Иван Опарин\\
Relux Works\\
\href{https://github.com/relux-works}{github.com/relux-works}\par}
\vspace{4pt}
{\normalsize Октябрь 2026, версия 7.0\par}
\end{center}

\vspace{6pt}

\begin{list}{}{\leftmargin=0.9cm\rightmargin=0.9cm}
\item\small
\textbf{Аннотация.} Языковая модель может предложить действие, но из предложения не должно следовать право его совершить. Swarma строит на этом разделении платформу для роёв автономных агентов, то есть групп агентов с разными ролями, которые вместе работают над целью. Платформа состоит из трёх слоёв. Swarma Runtime удерживает агентов на машинах одного владельца: у каждого агента своя учётка ОС, ключи хранит отдельная служба, каждое действие требует гранта, который при делегировании только сужается, а сообщения из мессенджера проходят детерминированные ворота доверия. Процесс работы подключается модулем: плейбук задаёт роли, маршрутизацию задач к исполнителям, политику ревью и правило приземления, а неизменные правила ядра доски запрещают принять изменение без независимой проверки точной ревизии. Swarma Protocol даёт независимым владельцам общую среду проекта с журналом под сертификатами кворума, оркестратором по аренде и продвижением основной ветки только к проверенному неизменяемому кандидату. Ключи, сообщения, учётки на носителе, гранты, память и аттестации сходятся на одной ячейке личности, которую Swarma Chain делает общей для всех проектов. Оценка эталонных компонентов: корпус подписанных записей из 2632 случаев проходит у генератора и двух независимых проверок, два вычислителя памяти согласны на 66 смысловых векторах, квалификационный набор на машине Linux проходит 11 из 11 проверок, ограниченная модель интеграции показывает необходимость четырёх правил, а эталонная проверка одной записи занимает от 2 до 18~мс в одном потоке.
\end{list}

\section{Введение}
\label{sec:intro}

Агент разработки читает файлы, вызывает shell и отправляет сообщения. Для операционной системы он остаётся обычным процессом пользователя, запустившего его, и получает все права этого пользователя: те же файлы, ту же конфигурацию, те же ключи. Косвенная инъекция в промпт \cite{greshake,agentdojo} пользуется именно этим. Допустим, вывод сборки предлагает ``для исправления'' открыть соседний файл с учётными данными. Модель принимает текст за подсказку и вызывает чтение. Если файл доступен процессу, исход зависит от того, сообразит ли модель отказаться. Если доступ закрыт границей ОС, ошибка модели файл не откроет.

Отсюда исходная позиция работы: модель может предложить действие, но предложение не даёт ей полномочия. Полномочие проверяет сторона, которая совершает эффект, детерминированным кодом и по проверяемым основаниям. Модель остаётся источником предложений, и её суждение входит в систему только там, где оно сужает допуск.

\textbf{Рой и его участники.} \emph{Агентом} мы называем участника работы с ролью, полномочиями и устойчивой криптографической личностью; его работу исполняет процесс агентской среды, например Claude Code или Codex. \emph{Роем} называются несколько агентов с разными ролями, которые вместе работают над целью. \emph{Владелец} отвечает за свои машины, агентов и оплату моделей, и ему принадлежит корень доверия установки. \emph{Оператор} означает человека владельца, который выдаёт полномочия и подтверждает действия, расширяющие права. \emph{Оркестратор} планирует работу роя: разбивает цель на задачи, назначает исполнителей и инициирует приёмку в пределах выданного ему полномочия. \emph{Продьюсер} производит изменение, а \emph{ревьювер} независимо проверяет его и выносит вердикт. Право принять результат в общий проект отделено от права его произвести и от права о нём судить.

Рой порождает вопросы, на которые модель не может ответить своим текстом. Кто отправил сообщение и что ему разрешено? Кто вправе дописать ввод в идущую сессию? Кто отвечает за работу, когда сессия остановилась на полуслове? К какому состоянию кода относится одобрение? Как принять ёмкость модели другого владельца, не отдав ему ключей проекта? Как агенту предъявить историю своей работы в новом проекте? Агентские фреймворки \cite{autogen,metagpt,chatdev,sweagent} координируют роли и инструменты, а эти вопросы оставляют среде исполнения.

\textbf{Зачем нужна каждая часть.} Swarma отвечает на эти вопросы отдельными механизмами, и у каждого механизма своя причина. Процесс агента получает права пользователя, поэтому каждый автономный агент работает на \emph{станции}, под своей учёткой ОС с сетью, ограниченной ядром. Ключ в доступном агенту файле можно украсть уговором, поэтому ключи держит отдельная служба, ключник, и агент только просит подпись. Подпись говорит, кто отправил байты, и ничего не разрешает, поэтому право на действие несёт грант, а ключи, гранты и отзывы ведёт реестр владельца. Членство в комнате мессенджера прав не даёт, поэтому мост проверяет каждое сообщение до доставки. У сессии много входов, поэтому вводом владеет хост сессий с единственным писателем. Знание и обязательства теряются вместе с контекстом, поэтому знание хранит память агентов, а обязательства со сроками ведёт Swarma Keeper. Принятие кода не имеет механического оракула, поэтому процесс работы задаёт подключаемый модуль с неизменными правилами ядра. Чужая ёмкость нужна без чужого доступа к среде, поэтому Swarma Protocol разделяет мышление, управление и исполнение. История участника должна пережить проект, поэтому все ключи сходятся на одной личности, а Swarma Chain делает её общей.

\textbf{Три слоя.} Swarma состоит из трёх слоёв (рис.~\ref{fig:layers}). \textbf{Swarma Runtime} держит рой на машинах одного владельца: учётки, ключи, полномочия, сообщения, сессии и окружение агентов. \textbf{Swarma Protocol} задаёт совместную работу независимых владельцев над единой средой исполнения проекта. \textbf{Swarma Chain} хранит общее состояние личностей, имён, аттестаций и расчётов за пределами отдельного проекта. Рантайм полезен сам по себе, протокол опирается на рантайм, цепочка связывает участников разных проектов.

Работа описывает платформу целиком. Границы доверия рантайма и протокола подробно разбирает работа-спутник \cite{swarmatrust}, устройство цепочки описывает отдельная работа \cite{swarmachain}, а связную модель с примерами и состояние разработки каждой части даёт книга \cite{swarmabook}. Где работа и спецификация модуля расходятся, права спецификация.

Вклад работы:
\begin{enumerate}
\item Модель системы и модель угроз для всей платформы, где захваченный инъекцией агент считается штатным противником на каждом слое (раздел~\ref{sec:model}).
\item Композиция рантайма, в которой допуск решает детерминированная проверка у исполнителя эффекта: станции с контрактом изоляции исполнения, ключник с определением собеседника ядром, гранты с ослаблением, реестр владельца со свидетелем отката, ворота доверия поверх Matrix, хост сессий с единственным писателем и память как потребитель реестра (раздел~\ref{sec:runtime}).
\item Рабочий процесс как подключаемый модуль: плейбук задаёт роли, маршрутизацию, политику ревью и правило приземления, а неизменные правила ядра действуют под любым плейбуком; процесс по умолчанию состоит из продьюсера, ревьювера, ищущего основания отклонить, и приземления точной проверенной головы (раздел~\ref{sec:workflow}).
\item Протокол независимых владельцев над одной средой: журнал под сертификатами кворума трёх уровней, оркестратор по аренде, принятие, адресованное точному кандидату, и федерация проектов (раздел~\ref{sec:protocol}).
\item Схождение ключей, сообщений, учёток, грантов, памяти и аттестаций на одной ячейке личности и её интеграция со Swarma Chain (разделы~\ref{sec:identity} и~\ref{sec:chain}).
\item Анализ гарантий по слоям и оценка эталонных компонентов: корпуса соответствия с независимыми реализациями, привилегированные запуски на машине Linux, ограниченная модель интеграции, микрозамеры проверки записей и статистика ревью (разделы~\ref{sec:analysis} и~\ref{sec:eval}).
\end{enumerate}

\begin{figure}[t]
\centering
\begin{tikzpicture}[yscale=0.9]
  \draw[zone] (-0.3,2.1) rectangle (4.6,4.6);
  \node[zl] at (-0.3,4.6) {Swarma Runtime владельца A};
  \draw[zone] (10.0,2.1) rectangle (14.9,4.6);
  \node[zl] at (10.0,4.6) {Swarma Runtime владельца B};
  \draw[zone] (5.1,1.2) rectangle (9.5,4.6);
  \node[zl] at (5.1,4.6) {Swarma Protocol};
  \draw[zone] (-0.3,-1.9) rectangle (14.9,0.5);
  \node[zl] at (-0.3,0.5) {Swarma Chain};
  \node[box] (ra) at (2.15,3.25) {станции, ключник,\\\scriptsize гранты, мост, хост сессий};
  \node[box] (rb) at (12.45,3.25) {станции, ключник,\\\scriptsize гранты, мост, хост сессий};
  \node[box] (env) at (7.3,3.55) {среда проекта\\\scriptsize основная ветка, области задач};
  \node[box] (core) at (7.3,2.05) {кор-ноды\\\scriptsize журнал, кворум};
  \node[box,minimum width=136mm] (chain) at (7.3,-0.75) {ячейки личностей и имена \quad пачки аттестаций \quad головы журналов \quad хэш блока \quad расчёты};
  \draw[arr] (ra) to node[lbl,above=1pt]{команды} (env);
  \draw[arr] (rb) to node[lbl,above=1pt]{команды} (env);
  \draw[arr] (core) to node[lbl,right=1pt]{QC} (env);
  \draw[arr] (ra.south) to node[lbl,left=2pt]{личности} (ra.south |- chain.north);
  \draw[arr] (rb.south) to node[lbl,right=2pt]{личности} (rb.south |- chain.north);
  \draw[arr] ([xshift=-6mm]core.south) to node[lbl,left=1pt]{аттестации} ([xshift=-6mm]core.south |- chain.north);
  \draw[arr] ([xshift=6mm]core.south |- chain.north) to node[lbl,right=1pt]{случайность} ([xshift=6mm]core.south);
\end{tikzpicture}
\caption{Три слоя Swarma. Рантайм каждого владельца удерживает своих агентов; протокол связывает рантаймы независимых владельцев с одной средой проекта, где команды исполняются только по записям журнала под сертификатом кворума (QC); цепочка хранит ячейки личностей, имена, пачки аттестаций и расчёты, общие для всех проектов.}
\label{fig:layers}
\end{figure}

\section{Модель системы и модель угроз}
\label{sec:model}

\textbf{Принципалы и роли.} Кроме владельца, оператора и агентских ролей из введения, в работе участвуют ещё четыре роли. \emph{Проект} задаёт цель, среду, правила приёмки и политику допуска участников. \emph{Контрибьютор} приносит проекту ёмкость своего агентского цикла, то есть модели с контекстом и вызовами инструментов, а \emph{донор} отдаёт квоту своей подписки на модель. \emph{Верификатор} повторяет критические проверки на другой машине. Совмещение ролей дополнительных прав не даёт: право дать ресурс, право владеть установкой и право изменить проект различаются.

\textbf{Хосты, станции и службы.} Хостом называется машина владельца, на которой работают агенты. \emph{Станцией} называется место одного агента на хосте: своя учётка ОС, своя личность на носителе сообщений, свои ключи, дом и корневой контекст; учётка с отдельным UID без остальных проверок станцией не считается. Привилегированные службы хоста разнесены по учёткам (рис.~\ref{fig:host}): ключник работает под root, менеджер пользователей и запускатель получают root на одну узкую операцию, остальные службы работают под своими служебными учётками.

\textbf{Носитель и узел доверия.} \emph{Носителем} называется мессенджер, который перевозит сообщения агентов и людей; носителем служит Matrix \cite{matrix}. Реестр владельца и свидетель отката работают на \emph{узле доверия}, отдельной машине вне административного домена платформы. Её снимают у провайдера, который от платформы не зависит, и администрирует её только владелец; у агентов там нет root, у платформы нет административного доступа.

\textbf{Протокол.} В Swarma Protocol \emph{кор-ноды} независимых владельцев ведут журнал полномочий и принятия, \emph{рабочими нодами} служат недоверенные агентские циклы контрибьюторов, а \emph{среда исполнения} проекта исполняет только разрешённые журналом действия.

\begin{figure}[t]
\centering
\begin{tikzpicture}[yscale=0.86]
  \draw[zone] (-0.2,0.9) rectangle (9.7,4.25);
  \node[zl] at (-0.2,4.25) {хост владельца};
  \draw[zone] (10.6,0.9) rectangle (14.9,4.25);
  \node[zl] at (10.6,4.25) {носитель};
  \draw[zone] (3.1,-2.05) rectangle (9.7,-0.25);
  \node[zl] at (3.1,-0.25) {узел доверия};
  \node[box] (a1) at (1.35,3.25) {станция A\\\scriptsize учётка, дом, сессия};
  \node[box] (a2) at (1.35,1.75) {станция B\\\scriptsize учётка, дом, сессия};
  \node[box] (kh) at (4.75,3.25) {ключник\\\scriptsize root};
  \node[box] (svc) at (4.75,1.75) {диспетчер, брокер,\\\scriptsize хост сессий};
  \node[box] (hl) at (8.15,3.25) {помощник допуска\\\scriptsize ключ машины};
  \node[box] (br) at (8.15,1.75) {мост\\\scriptsize служебная учётка};
  \node[box] (rg) at (12.75,3.25) {регистратор};
  \node[box] (cs) at (12.75,1.75) {сервер Matrix};
  \node[box] (reg) at (4.75,-1.2) {реестр владельца};
  \node[box] (wit) at (8.15,-1.2) {свидетель отката};
  \node[box] (ph) at (0.95,-1.2) {телефон оператора\\\scriptsize ключ в защищённом чипе};
  \draw[arr] (a1) to node[lbl,above=1pt]{UID} (kh);
  \draw[arr] (a2) to (svc);
  \draw[arr] (hl) to node[lbl,above=1pt]{запрос} (rg);
  \draw[arr] (rg) to node[lbl,right=1pt]{пропуск} (cs);
  \draw[arr] (cs) to node[lbl,above=1pt]{события} (br);
  \draw[arr] (reg) to node[lbl,right=1pt]{снимки $\le$60~с} (svc);
  \draw[arr] (reg) to (wit);
  \draw[arr] (ph) to (reg);
\end{tikzpicture}
\caption{Хост владельца. Агенты на станциях ключей не держат; ключник и помощник допуска узнают собеседника по UID от ядра. Регистратор выдаёт пропуск на носитель, мост проверяет сообщения, реестр на узле доверия принимает подписанные действия оператора и выдаёт службам снимки, свидетель фиксирует его контрольные точки.}
\label{fig:host}
\end{figure}

\textbf{Противники.} У трёх слоёв разные противники, и одна мысль на каждом слое требует своего механизма. Мы рассматриваем восемь противников:
\begin{enumerate}[label=П\arabic*.]
\item Захваченный агент подчинился внедрённому тексту и пользуется всеми возможностями своего процесса, включая подписи, которые ключник выдаёт его личности.
\item Чужой участник, то есть агент или человек другого владельца, рабочая нода протокола или гость комнаты, пытается расширить полномочия, подменить объект одобрения или действовать от чужого имени.
\item Сетевой противник подменяет, задерживает и повторяет сообщения, а сервер носителя может выпустить событие от имени любого адреса своего домена.
\item Противник получил root на одной зарегистрированной машине из нескольких; корни владельца вне этой машины, ключ оператора и узел доверия ему недоступны.
\item Противник на узле доверия похитил ключ реестра или пытается откатить и раздвоить журнал реестра.
\item Злонамеренный оркестратор протокола предлагает вредные задачи и пытается сам выбрать проверяющего или интегрировать результат.
\item До $f$ кор-нод злонамеренны на уровне доверия L2, меньшинство ключей скомпрометировано на L1.
\item Криптографический противник ломает класс подписей на эллиптических кривых.
\end{enumerate}

\textbf{Допущения.} Для гарантий хоста доверены его ядро, root и доверенная база платформы: ключник, менеджер пользователей, запускатель, диспетчер, брокер, хост сессий, мост с воротами доверия и помощник допуска, на узле доверия реестр, его процесс подписи и свидетель, вне хоста регистратор и телефон оператора. Мастер-секрет ключника цел. Оператор подтверждает расширяющие действия на телефоне с ключом в Secure Enclave или StrongBox и видит подписываемое содержание. Владелец честен перед своими агентами и единолично администрирует узел доверия. У установки есть проверенный или контролируемый оператором источник времени, например NTS. В протоколе среда проекта доверена в отношении целостности, а провайдер модели доверен в отношении конфиденциальности в пределах условий, которые проект объявил до присоединения.

\textbf{Вне рамок.} Платформа не защищает от эксплойта ядра и не сохраняет гарантий хоста с захваченным root (радиус поражения для остальных машин разобран как П4). Она не обещает правильности суждения модели, конфиденциальности кода от контрибьютора и его провайдера и полноты наблюдения за недоверенной нодой вне протокола. Побочные каналы и подделка оборудования не рассматриваются. Отказ в обслуживании не предотвращается: при недоступном свидетельстве система закрывает допуск и сознательно жертвует живучестью ради целостности.

\textbf{Задача.} Ни одно действие агента с эффектом за пределами его станции не исполняется без проверяемого основания: подлинного отправителя, актуального гранта, одобрения, адресованного точному объекту, и физического удержания исполнения в выданной области. Основание проверяет сторона, совершающая эффект. Отсутствие свидетельства, подтверждённое отсутствие события и невозможность прочитать состояние различаются, и разрешения ни одно из них не даёт.

\section{Swarma Runtime на машинах одного владельца}
\label{sec:runtime}

Swarma Runtime собран из небольших служб с именованными контрактами под разными пользователями ОС; границу задают учётки, права файлов, сокеты и правила ядра.

\subsection{Станция и контракт изоляции исполнения}

Каждый автономный агент работает отдельным пользователем ОС: оркестратор, постоянный специалист, временный продьюсер и рабочее место удалённого участника. Собственный UID отделяет файлы и процессы одного агента от другого, и процесс продьюсера не прочитает рабочий каталог ревьювера, даже зная путь. Контейнеры и виртуальные машины для этого не нужны: рантайму нужна граница между агентами одной машины, а отдельного образа окружения на агента ему не требуется. На Linux у каждого поколения учётки своя cgroup v2 под управлением root с лимитами \texttt{cpu.max}, \texttt{memory.max}, \texttt{memory.swap.max}, \texttt{pids.max} и лимитами ввода-вывода, а также квоты файловой системы. На macOS роль этой границы играет квалифицированный профиль Seatbelt для класса задачи. Агент считается остановленным, когда у всей его учётки нет процессов в течение опубликованного периода тишины; выход ведущего процесса этого не доказывает.

Привилегированных входов два. Менеджер пользователей (\texttt{swarma-user-manager}) создаёт и выводит учётки и управляет их сетевыми правилами и группами ресурсов, запускатель (\texttt{swarma-user-launcher}) запускает одну одобренную программу под активной учёткой. Обе программы вызываются через \texttt{sudo} без аргументов и принимают один JSON ограниченного размера, а план запуска приходит неизменяемыми байтами на дескрипторе 3. Каждое создание учётки получает случайное 128-битное \emph{поколение}, и на поколение опираются привязки, гранты и группа ресурсов, поэтому повторно выданный UID чужих прав не наследует. Если ответ о шаге потерян, диспетчер наблюдает тот же эффект, и отсутствие свидетельства не позволяет считать, что эффекта не было.

Сеть ограничивает ядро: правила \texttt{pf} на macOS или \texttt{nftables} на Linux пускают агента только к слушателю доверенного прокси, своему у каждого поколения учётки. Прокси пропускает \texttt{CONNECT} к адресатам провайдера и ограничивает имя, порт и проверенный адрес; TLS внутри туннеля он не расшифровывает, поэтому разрешённый адресат остаётся каналом выноса данных. У рабочего места удалённого участника исходящей сети нет совсем, а сборки и тесты кандидата идут на эфемерных раннерах другой машины без секретов проекта. Публикацию кандидата, обновление базы и приземление выполняет отдельная служба действий с репозиторием под своей учёткой, по подписанному доской запросу на точного кандидата; голову ветки она продвигает сравнением-и-заменой.

Что станция гарантирует, задаёт \emph{контракт изоляции исполнения}: профиль изоляции, вектор возможностей, привязки личности, размещения и исполнения, квитанция границы и кортеж квалификации. Запускатель применяет допущенную границу и пишет защищённую квитанцию, диспетчер сверяет покрытие, брокер перед выдачей учётных данных проверяет текущую квитанцию. Гарантии названы 42 идентификаторами общего словаря; например, \texttt{isolation.read} означает, что работа станции не читает чужих личных данных и секретов платформы. У каждой гарантии на каждой платформе один из четырёх статусов: \texttt{present}, \texttt{weaker-with-compensation}, \texttt{absent}, \texttt{PENDING-VERIFY}. Задача объявляет нужные гарантии записью \texttt{station-required-guarantees/1}; \texttt{absent} и \texttt{PENDING-VERIFY} дают отказ, а ослабление допустимо, только если оператор отдельным аутентифицированным действием принял его для этой задачи, станции, компенсации и срока. Машину владелец квалифицирует сам воспроизводимой процедурой; аттестацию подписывает отдельный ключ машины и называет каждую из 42 гарантий ровно один раз. При дрейфе защищённых входов или проваленной проверке новые старты закрываются сразу, а работающее исполнение сужается или останавливается целиком не позже 300~с.

\subsection{Ключник: личности без ключей у агента}

\emph{Ключник} (Swarma Keyholder) служит привилегированной службой с мастер-секретом установки, которая выполняет криптографические операции за агентов. Закрытого ключа личности у агента нет: инъекция действует руками агента, и агента с ключом в доступном ему файле достаточно уговорить этот файл прочитать. Где возможно, мастер-секрет держит аппаратное хранилище. Ключи агентов выводятся жёсткой деривацией по BIP-32 в варианте SLIP-0010 для Ed25519 \cite{bip32,slip10}, а путь \texttt{назначение'/\allowbreak продукт'/\allowbreak проект'/\allowbreak агент'} кодирует стабильную личность, отдельную от UID и пути в файловой системе. Личность поэтому переживает пересоздание учётки и переезд на другую машину с тем же мастер-секретом.

Ключник слушает Unix-сокет и получает UID собеседника от ядра (\texttt{SO\_PEERCRED} на Linux, \texttt{LOCAL\_PEERCRED} на macOS). Личность выбирается по привязке с UID, поколением учётки, путём, владельцем и номером регистрации, и подключение допустимо, только если поколение этого UID активно в журнале менеджера пользователей. Процесс ревьювера может написать в запросе имя оркестратора, но его соединение всё равно несёт UID ревьювера. Сообщения агента ключник подписывает отдельным ключом сообщений Ed25519 с путём \texttt{m/4'/\allowbreak продукт'/\allowbreak проект'/\allowbreak агент'}, и только документ с меткой \texttt{swarma-msg/2} с прообразом не длиннее 24\,576 байт. Принадлежность ключа сообщений агенту удостоверяет запись \texttt{attestation/1}, подписанная корнем аттестаций ключника, сроком не длиннее 86\,400~с; подпись открыта, только если аттестация действует и опубликована в реестре владельца.

Каждый корень подписывает только свой язык записей (аттестации ключника, действия оператора, гранты, контрольные точки реестра, регистрации машин), и совпадение ключей двух ролей ведёт к отказу. Выдача гранта, подключение подписки и другие расширяющие действия требуют человека: приложение на телефоне оператора показывает точное содержание действия и после подтверждения биометрией подписывает именно его ключом в защищённом чипе, а второе устройство обязательно. Общий стоп всей установки выполняется одной подписанной командой, которая доставляется напрямую каждой службе; возобновить работу можно только подписью оператора.

\subsection{Гранты, реестр владельца и свидетель}

\emph{Грантом} называется подписанное утверждение: кто выдал полномочие, кому, на какие действия и цели, при каких условиях, до какого срока и можно ли передать его дальше. Делегирование только сужает: дочерний грант покрывает подмножество действий и целей родителя, его срок не длиннее, глубина в разрешённых пределах. Цепочку до корня проверяет исполнитель действия, даже если раньше на маршруте её уже проверили. В общем языке грантов ключ агента бывает субъектом и листом цепочки, но издателем звена не бывает, поэтому захваченный оркестратор не раздаёт полученное соседям; контракт брокера v0 допускает одно исключение, суженный грант, который родительский оркестратор подписывает создаваемому им оркестратору. Вызывающий остаётся в цепочке потомка, и в контракте брокера учётных данных v0 лист агента равен пересечению запроса и двух цепочек, вызывающего и диспетчера, с одним действием \texttt{credential.lease}, глубиной 0 и не больше четырёх конвертов в цепочке. Если цепочку построить нельзя, ответ \texttt{grant\_unconstructible} приходит до любых эффектов. Отзыв гранта оркестратора поэтому отрезает его агентов при следующей аренде или продлении.

\emph{Реестр владельца} служит единственным авторитетным писателем ключей, аттестаций, грантов, отзывов и журналов личностей. Он хранит, упорядочивает и публикует записи, а содержимое не переподписывает. Собственный ключ реестра подписывает только снимки отзывов своей области, контрольные точки и записи свежести, поэтому кража этого ключа не даёт выпустить грант или аттестацию. Журнал реестра состоит из записей \texttt{journal-entry/1} со ссылкой на предыдущую, а подписанная \texttt{checkpoint/1} называет голову и текущий снимок отзывов. Рядом работает \emph{свидетель отката}: служба под отдельной учёткой и процессом, которая по запросу с вызовом называет последнюю голову и хранит доверенную контрольную точку вне узла. Две разные контрольные точки с одним номером дают \texttt{journal\_fork}, точка с номером меньше засвидетельствованного даёт \texttt{journal\_rollback}.

Потребители принимают снимок отзывов не старше 60~с; возраст неизменного снимка продлевает запись свежести, которую реестр публикует каждые 20~с, вместе со свежим чтением головы у свидетеля. Записи отзыва накапливаются. Пусть $S$ означает время выпуска первого подписанного полного снимка с записью отзыва, $A\le60$~с допустимый возраст снимка, $F$ время вступления записи в силу. Тогда запись действует на каждой проверке не позже $\max(S+A,F)$ и никогда раньше $F$. Проверка отзыва даёт три исхода: не отозвано, отозвано, недоступно; два последних закрывают допуск. Время проверяется по водяному знаку, и при потере проверенного времени службы держатся не дольше 10 минут по монотонным часам. Сдерживание шире отзыва гранта: одна операция ``отозвать и остановить'' отзывает полномочие брокера и отгораживает исполнение и сеть, а о токене провайдера, который уже лежит в процессе агента, честно сообщает, что он пригоден до аннулирования у провайдера.

\subsection{Подписанные сообщения и обмен агентов поверх Matrix}

Сообщение агенту может быть просьбой коллеги, командой владельца или чужим текстом, выданным за распоряжение, и для модели все три выглядят одинаково. Поэтому доставка отделена от полномочий. Носитель перевозит событие, мост проверяет происхождение и допустимость, хост сессий решает, когда передать допущенное событие среде агента. У агента нет учётных данных мессенджера, и прямой выход к носителю ему закрывает файрвол.

Учётку на носителе агент получает без общего секрета. Помощник с ключом машины узнаёт UID вызывающего от ядра и подписывает только запрос на допуск с именем, разрешённым этому UID; \emph{регистратор} (\texttt{swarma-registrar}) сверяет ключ машины с реестром машин в git и выдаёт пропуск, JWT с подписью Ed25519 на 120~с. Сервер Matrix доверяет только ключу регистратора и выдаёт токен доступа на 86\,400~с без refresh-токена, поэтому скомпрометированная машина заводит учётки только со своей меткой.

Подписанную запись платформы мы называем \emph{конвертом}: он состоит из канонической формы байтов с доменной меткой и подписи автора над ними, а формат и порядок проверок задаёт спецификация подписанных записей \texttt{swarm-envelope/2}. \emph{Мост} стоит между носителем и сессией агента и работает отдельным экземпляром на каждую личность станции под общей служебной учёткой сообщений, никогда под учёткой агента. Провайдер Matrix сначала дописывает событие в почтовый файл и сбрасывает его на диск, затем продвигает курсор синхронизации, а допуск идемпотентен по паре отправителя и идентификатора события. \emph{Ворота доверия} проверяют сообщение детерминированно, без сети и без вопроса к модели, в строгом порядке: каноническая форма байтов, ключ в снимке реестра, подпись, цепочка грантов до корня с сужением действий и целей, снимок отзыва, окно времени, журнал повторов и ревизия привязки цели. Затем небольшая локальная модель (decision guard) может только сузить допуск, и её недоступность даёт отказ. Допущенный текст во ввод сессии не вставляется: хост сессий получает короткий звонок, а тело и запись проверки агент читает из почтового ящика как данные. Сообщения других агентов остаются данными при любой подписи, и пересказ решения владельца другим агентом полномочий автора не переносит.

Сила свидетельства об отправителе-человеке задаётся лестницей. На уровне A0 отправителя удостоверяет сервер носителя владельца, на A1 устройство отправителя подписано закреплённым мастер-ключом перекрёстной подписи Matrix, на A2 ключ человека подписал точные байты конверта и ворота проверили грант, на A3 оператор подтвердил точное действие на телефоне. Расширяющие действия требуют A3. Исходящий путь сохраняет точные байты: агент просит ключник подписать текст, а мост проверяет подпись и допустимость аудитории комнаты и не переписывает его. У каждой станции есть личная комната с владельцем ровно на двух участников; владельца записывает поднимающая станцию сторона вне области записи агента.

\subsection{Хост сессий и запуск}

Запуск разделён между программами, чтобы сборщик плана не получал привилегий запускателя. \emph{Swarma Curator} (далее Curator) управляет окружением агентов: собирает из профиля инструкции, скиллы, подключения MCP и настройки и выдаёт для каждой агентской среды управляемый дом. План запуска фиксирует профиль с пином, то есть хэшем, дом, сетевой профиль и режим учётных данных; собирает его вызывающий, а исполнитель под учёткой агента только проверяет байты. Любой MCP подключается с явного согласия, а неизвестная готовность логина запуска не разрешает. Диспетчер (\texttt{swarma-dispatcher}) проверяет подпись заявки, цепочку грантов, свежий вердикт отзыва, квоту и ёмкость машины, а брокер учётных данных (\texttt{swarma-credential-broker}) выдаёт аренду токена провайдера только стороне агента и при каждом продлении повторяет проверки.

\emph{Хост сессий} управляет жизнью долгой сессии и состоит из службы \texttt{swarma-session-host} под своей учёткой и раннера \texttt{swarma-session-runner} под учёткой агента, который держит PTY или app-server среды. Главный инвариант: во ввод сессии пишет один диспетчер ввода хоста, и только допущенные элементы долговечной очереди. Раннер принимает команды только текущей эпохи контроллера, поэтому соединение, пережившее сбой службы, права старого контроллера не сохраняет. Элемент считается применённым только по подтверждению среды; запись байтов в канал применения не означает. Забор допуска разрешает или запрещает каждый эффект сессии, а контролируемый перезапуск закрывает доставку, дожидается покоя, сохраняет пакет передачи и подтверждает отсутствие прежнего исполнения, прежде чем запустить преемника.

Цель сессии не принадлежит агенту. Её ставит вызывающий при запуске или публикует доска, и у учётки самого агента прав на операции с целями нет; сообщение агента о выполнении цели остаётся предложением исхода. Сессия служит кэшем: авторитетны журнал хоста, записка о передаче и, если есть доска, её задача и цель. После перезагрузки постоянные команды поднимает Swarma Keeper командой \texttt{keeper up} через диспетчер, а остановленное человеком остаётся остановленным.

\subsection{Память агентов}

Память агентов хранит знание отдельно от сменяющихся исполнителей: типизированные утверждения вместе с источниками, обоснованиями, оценками и ревизиями. На запрос она отвечает, какие допущенные свидетельства поддерживают или опровергают точно очерченное утверждение, с чьей точки зрения и на какой момент; пустой ответ не доказывает ложности утверждения. Агент получает память как данные с происхождением каждого элемента, и текст из памяти инструкцией не становится.

С реестром владельца память стоит на одной подложке подписанных записей и читает его как потребитель: отправителя записи определяет по ячейке личности в текущем снимке, основание представительства берёт из гранта, а недоступный реестр никогда не считается выданным грантом. Полномочий память не выдаёт: у неё нет операций одобрить, принять задачу, выдать грант, запустить или разбудить. В профиле git-пространств агент сам подписывает запись через ключник и отправляет её, а каждый читатель проверяет подпись, схему и право на ревизию при чтении. Содержимое этого профиля ограничено нечувствительным, потому что из истории git ничего не стирается; чувствительные данные обслуживает защищённая служба с единственным писателем. Нормативный пакет 0.3.0 содержит 53 пункта ядра, 14 закрытых схем и 85 векторов соответствия: 77 смысловых и 8 для защищённого развёртывания.

\section{Рабочие процессы как подключаемый модуль}
\label{sec:workflow}

Механизмы рантайма говорят, кто вправе действовать и где исполняется действие. Они не говорят, как команда устраивает работу: какие роли есть, кто производит, кто проверяет и по какому правилу результат попадает в основную ветку. У разных команд ответы разные, а гарантии рантайма от этих ответов зависеть не должны. Поэтому процесс работы оформлен подключаемым модулем.

\subsection{Модуль процесса и его состав}

\emph{Модулем процесса} называется пакет, который оркестратор загружает, чтобы вести работу: скилл оркестратора, скиллы ролей, поставщики инструментов и \emph{плейбук}, декларативное описание процесса в файле \texttt{Playbook.json}. Плейбук задаёт типы работ, состояния, владельцев состояний и разрешённые переходы с проверяемыми условиями (guard), наборы проверок и делегирование действий с потолками полномочий. Плейбук содержит только данные, без исполняемых программ, скриптов и URL; инструмент предоставляет установленный скилл, и запускается инструмент, только если проект его включил. Пакет ставит Curator, а скиллы ролей попадают только в запуски своих ролей: если бы их раскладывали в проект, любая сессия нашла бы их, и оркестратор начал бы действовать по инструкциям ревьювера.

Модуль процесса отвечает на четыре вопроса команды. Первый: какие роли существуют и кто владеет каждым состоянием; владельцем может быть агент со скиллом роли, человек из группы, инструмент или другой процесс, и в каждый момент у элемента ровно один владелец. Второй: как задача доходит до исполнителя; маршрутизатор ролей превращает владельца нового состояния в запрос запуска (раздел~\ref{sec:routing}). Третий: какая проверка нужна для принятия, то есть политика ревью. Четвёртый: по какому правилу результат приземляется в основную ветку. Описание установленного, работы и окружения доски разведено по трём файлам: \texttt{Skillfile.json} читает Curator, \texttt{Playbook.json} описывает работу, \texttt{task-board.config.json} задаёт окружение доски.

Плейбуки компонуются: проект может взять опубликованный процесс и описать отличия либо соединить несколько процессов через слоты, например Scrum для выбора работы и сроков с циклом разработки для пути конкретного изменения, а в iOS-проекте заменить скилл разработчика и вставить состояние QA. Опубликованные части закрепляются хэшами. Исполняется версия, вступившая в силу, а новая требует проверки и подтверждения человеком; отчёт о полномочиях выделяет изменения, которые уводят работу мимо человека или проверки.

\begin{figure}[t]
\centering
\begin{tikzpicture}[yscale=0.9]
  \node[box] (pb) at (0,3.2) {модуль процесса\\\scriptsize плейбук, скиллы ролей};
  \node[box] (core) at (4.4,3.2) {ядро доски\\\scriptsize неизменные правила};
  \node[box] (rr) at (8.8,3.2) {маршрутизатор ролей\\\scriptsize и маршрутизатор моделей};
  \node[box] (disp) at (13.0,3.9) {диспетчер\\\scriptsize разовый запуск};
  \node[box] (wk) at (13.0,2.4) {хост сессий\\\scriptsize верстак};
  \node[box] (prod) at (13.0,0.6) {продьюсер\\\scriptsize своя станция};
  \node[box] (cr) at (8.8,0.6) {Change Request\\\scriptsize неизменная ревизия};
  \node[box] (rev) at (4.4,0.6) {ревьювер\\\scriptsize другой издатель модели};
  \node[box] (del) at (0,0.6) {служба действий\\\scriptsize с репозиторием};
  \node[box] (main) at (0,-1.1) {основная ветка};
  \node[box] (kp) at (6.6,-1.1) {Swarma Keeper\\\scriptsize обязательства и сроки};
  \draw[arr] (pb) to node[lbl,above=1pt]{версия в силе} (core);
  \draw[arr] (core) to node[lbl,above=1pt]{переход} (rr);
  \draw[arr] (rr) to (disp);
  \draw[arr] (rr) to node[lbl,below=2pt]{цель} (wk);
  \draw[arr] (wk) to (prod);
  \draw[arr] (prod) to node[lbl,above=1pt]{ревизия} (cr);
  \draw[arr] (cr) to node[lbl,above=1pt]{инертный diff} (rev);
  \draw[arr] (rev) to node[lbl,above=1pt]{вердикт} (del);
  \draw[arr] (del) to node[lbl,right=1pt]{CAS точной головы} (main);
  \draw[arr] (kp) to node[lbl,above=1pt]{передача} (cr);
\end{tikzpicture}
\caption{Процесс по умолчанию в модуле процесса. Ядро доски проверяет каждый переход по действующей версии плейбука, маршрутизатор ролей передаёт работу разовому запуску или постоянному верстаку, ревьювер получает инертный diff неизменной ревизии, а основную ветку продвигает только служба действий с репозиторием сравнением-и-заменой (CAS) точной проверенной головы. Swarma Keeper ведёт обязательства между участниками.}
\label{fig:workflow}
\end{figure}

\subsection{Доска и правила ядра под любым плейбуком}

\emph{Доска задач} (\texttt{task-board}) хранит работу, её состояние и свидетельства приёмки и исполняет плейбук. Элементы образуют иерархию Epic, Story, Task или Bug с описанием, прогрессом, статусами, ресурсами и историей; агент читает и меняет доску языком запросов и мутаций, контракт которого можно спросить у самой программы. В рантайме доска служит необязательным клиентом: она просит запуск у диспетчера и публикует цели через хост сессий, а платформенные модули от неё не зависят. Хранилище доски границей авторизации не служит: роль связана с грантом, и переход проверяет ядро доски.

Настраиваемый процесс может сужать разрешённое, но не может отменить предикаты ядра. Фиксированный слой скомпилирован в программу доски, и ни автор плейбука, ни человек, подтвердивший его версию, права отключить его не получают. Защищённое состояние определяется классом, поэтому переименование интеграции обхода не открывает. Неизменные условия таковы: задача с Change Request доходит до интеграции только с принятым запросом; работа идёт в дереве своей Story, и устаревшее дерево не интегрируется; история продвигается подписанной перемоткой точной головы с доказательством стыка; терминальный исход ревью требует свидетельства для этой ревизии; мутация и запись аудита пишутся одной транзакцией в журнал, который только дописывается. Отсюда простой запрет: нельзя написать плейбук, в котором изменение принято по одному завершению продьюсера. Дополнительные ревью, состояния QA и ручное согласование допустимы, удалять обязательное свидетельство нельзя.

Свидетельство перехода должно быть свежим: задача, вернувшаяся из ревью в разработку и снова пришедшая на ревью, требует вердикта текущего прохода. Инструменты проверки запускает исполнитель ядра в очищенном окружении, иначе агент, запросивший переход, мог бы подложить свою программу проверки через \texttt{PATH}.

\subsection{Маршрутизация к исполнителям: верстаки, разовые запуски и модели}
\label{sec:routing}

Ядро записывает состояние и намерение, а маршрутизатор ролей превращает владельца нового состояния в запрос запуска: профиль Curator с точным закреплением, класс задачи, привязки ввода-вывода и нужный контекст. Для разовой работы запрос уходит диспетчеру и проходит транзакцию запуска из раздела~\ref{sec:runtime}. Для постоянной сессии, если на неё есть право, доска публикует новую ревизию цели, и хост сессий доставляет её на границе хода.

Постоянным исполнителем служит \emph{верстак} (Werkstatt), постоянное рабочее место агента, то есть связка ``проект, роль, среда''. У верстака свой управляемый дом Curator с профилем роли, свой клон репозитория, своя личность на носителе и запись в реестре верстаков на доске с описанием запуска: среда, модель и усилие рассуждений, профиль, клон и личность. Поднимает и останавливает верстак хост сессий; без него сессию нельзя было бы разбудить, проверить или восстановить. Задачи приходят верстаку сообщениями в комнатах проекта со звонком, но право на задачу даёт запись на доске, а текст сообщения права не даёт. Разработчик сдаёт работу ревизией, ревьювер пишет вердикт на доску и в репозиторий не пишет. Клоны верстаков лежат в отдельном корне, выше которого нет файлов \texttt{CLAUDE.md} и \texttt{AGENTS.md}: агентская среда ищет их от рабочей папки вверх, и подмена дома этот подъём не останавливает.

Модель и среду выбирает \emph{маршрутизатор моделей} (\texttt{cmr}). Перед запуском разные владельцы отвечают на четыре вопроса: разрешено ли (доска по политике оператора), совместимо ли (реестр моделей), подходит ли (маршрутизатор) и доступно ли сейчас (предварительная проверка доски). Маршрутизатор ранжирует только внутри допущенного набора и текста задачи не читает: он получает роль, класс работы, сложность, режим бюджета и закрепки. Выбор детерминирован, без вызова модели, а запись решения с идентификатором SHA-256 канонического JSON позволяет пересчитать его позже. Для процесса существенно одно правило: роль ревьювера может требовать другого издателя модели, чем у запуска, который произвёл работу. Независимость проверяется по записанному происхождению того запуска, и дополнительное ревью того же издателя нарушенного требования не восстанавливает \cite{selfpref,llmjudge}. Тот же маршрутизатор выбирает верстак под задачу: сначала фильтры по роли, проекту, живости, доверию и независимости, потом качество, загрузка и цена; если подходящего верстака нет, он отказывает.

\subsection{Процесс по умолчанию: продьюсер, ревьювер и приземление точной головы}

Процесс по умолчанию состоит из трёх ролей и одного правила. Продьюсер работает на своей станции в самостоятельном репозитории со своим каталогом \texttt{.git}; связанный \texttt{git worktree} между разными учётками не годится, потому что через общие служебные данные тянет настройки и хуки родителя. Продьюсер публикует \emph{Change Request}, запрос на изменение с конкретной неизменяемой ревизией, и вердикт относится именно к ней. Ревьювер получает инертные diff и blob: только для чтения, с дайджестами, вне каталогов, где агентская среда ищет настройки проекта. Инструкции, хуки, MCP и разрешения кандидата не подхватываются ни при запуске, ни при вложенном чтении, возобновлении или входе дочернего агента.

Ревьювер в процессе по умолчанию ищет основания отклонить кандидата. Он повторяет проверки, сверяет изменение с контрактом и выносит вердикт с находками по классам серьёзности; один запуск не может принять собственную ревизию, и заказчик раунда свой результат не проверяет. Изменения в CI, хуках, инструкциях агентов и ссылках одобряет детерминированный предикат доски над текущим контрактом, политикой риска, точным кандидатом и полным множеством помеченных путей, вычисленным по накопленной дельте от базы до кандидата. Поэтому инструкция агенту, добавленная в первом непринятом кандидате, не пропадает из отчёта второго только потому, что второй правит лишь текст ошибки.

Правило приземления: основную ветку продвигает служба действий с репозиторием, а у агентов и у процесса доски такого права нет. Служба выполняет подписанную перемотку точной проверенной головы по подписанному доской запросу. Если основная ветка ушла вперёд, возникает другой кандидат, вердикт по сходству diff на него не переносится, и нужны обновление базы и повторное ревью. У процедуры есть предел. Пусть десять процентов агентов злонамеренны, а честный ревьювер находит подсаженную уязвимость в семидесяти процентах случаев; тогда при независимых ошибках один ревьювер пропускает её с вероятностью 0,37, два с вероятностью около 0,14, три около 0,05. Общая модель и одинаковый контекст коррелируют ошибки, поэтому расчёт иллюстрирует допущения и качество Swarma не измеряет; измерять нужно обнаружение дефектов на подходящих задачах (раздел~\ref{sec:eval}).

\subsection{Swarma Keeper: обязательства между участниками}

Процесс задаёт, кто владеет состоянием; он не отвечает, кто кому что должен прямо сейчас, когда работа переходит между участниками. Эту задачу решает \emph{Swarma Keeper}, кипер, главный цикл узла. Кипер следит, чтобы у каждого обязательства был держатель; путать его с ключником нельзя, потому что управление обязательствами и владение ключами имеют разные границы доверия. У кипера две части. Роль исполняет модель, которая разбирает просьбы и принимает решения в своих полномочиях, а модуль, команда \texttt{keeper} с библиотекой, ведёт журнал, проверяет переходы и сверяет сроки. Когда модель остановилась или её контекст сжали, часы обязательств не останавливаются.

Обязательство значит ``участник должен другому участнику сделать определённую работу к сроку''; виды обязательств: ответить, нести обязанность, доставить, проверить и решить. Передача проходит три этапа: отправлена, принята, начата. Отправитель называет срок принятия \texttt{accept\_by} и срок первого действия \texttt{act\_by}, ответ ``беру'' подтверждает только принятие, и до проверенного первого действия нового держателя обязанность остаётся у прежнего. У каждой попытки свой неповторяющийся \texttt{handover\_id}, а срок принимается только полной датой и временем UTC, поэтому он восстанавливается после сжатия контекста без догадки о дне.

Эскалация идёт вверх по четырём ступеням: оркестратор держателя, кипер продукта, мета-кипер владельца, человек-владелец. На шаг приходится один переход и одно уведомление, продлений не больше двух на эпизод, недоступного получателя пропускают с записью, а дойдя до человека, подъём останавливается, обязательство остаётся открытым, и молчание согласием не считается. Если автоматическое пробуждение не доказано, каждый шаг приходит двумя уведомлениями: адресным вопросом и видимым человеку \texttt{owner-needed}. Журнал дописывает один процесс кипера под его системной учёткой, записи он подписывает своим ключом через ключник, а состояние получается свёрткой записей по порядковому номеру, без опоры на настенные часы. Доску кипер не заменяет: непройденное ревью остаётся фактом доски, а журнал ссылается на него, если из него следует обязанность одного участника перед другим.

\subsection{Другие процессы под тем же ядром}

Смена процесса не создаёт второго механизма авторизации: у дизайна и разработки разные состояния, но владелец в обоих случаях предлагает переход ядру. Процессы вне разработки пользуются той же идеей точного предмета: если заказ дороже 500 евро идёт менеджеру, одобрение заказа на 620 евро не разрешает исполнить заказ на 6\,200, потому что одобрение связано с отпечатком полей. Проектировщик процессов, системная роль вокруг кипера, начинает с простейшего работающего процесса и отправляет изменения плейбука человеку на подтверждение, а валидатор процесса проверяет достижимость состояний, наличие владельцев, обходы guard и детерминизм.

\section{Swarma Protocol и федерация}
\label{sec:protocol}

\subsection{Одна среда и три плоскости}

Средой исполнения называется место, где лежит рабочее дерево проекта, выполняются команды и проверки и продвигается основная ветка. Проект может зависеть от Xcode, лицензий, устройств, GPU, ключей подписи или данных, которые нельзя вывозить, и копия репозитория копии такой среды не даёт. Протокол оставляет среду в одном авторитетном экземпляре и разделяет три плоскости. \emph{Мышление} выполняют агентские циклы контрибьюторов с их моделями и учётными данными. \emph{Управление} выполняют кор-ноды, которые ведут журнал полномочий и принятия. \emph{Исполнение} выполняет пассивная среда: агентских циклов она не запускает и исполняет только разрешённые команды. По сети идут команды, вывод, хэши и подписи; учётные данные контрибьютора остаются у него, а секреты проекта есть только у доверенного раннера с узкими именованными операциями.

\subsection{Журнал, кворум и уровни доверия}

Журнал представляет собой упорядоченную цепочку подписанных записей с $h_i = H(h_{i-1}\,\|\,c(r_i))$, где $c(r_i)$ означает каноническую форму записи. Запись совершена, когда у неё есть сертификат кворума (QC), набор подписей под одним содержимым, достаточный для уровня доверия проекта; до этого предложение записи среду не авторизует. Уровень выбирает проект:
\begin{itemize}
\item на L0 одна кор-нода даёт QC из одной подписи, а кор и среда доверены;
\item на L1 при $n\ge2$ порог $t>n/2$ защищает от одиночного несанкционированного решения при компрометации меньшинства ключей, а единственность истории держится на честном секвенсере;
\item на L2 при $n\ge3f+1$ византийски устойчивый консенсус \cite{pbft,hotstuff} использует порог $q=\lceil (n+f+1)/2\rceil$ \cite{byzquorum}, при $n=3f+1$ равный $2f+1$.
\end{itemize}
Фиксированные $2f+1$ при большем $n$ неверны: при $n=5$ и $f=1$ они дали бы порог 3, а два кворума по три подписи из пяти могут пересечься ровно в одном злонамеренном узле. Среда работает как лёгкий клиент: ей нужны подтверждённое назначение и доказательство его включения в дерево Меркла журнала, а наивысшую принятую версию журнала она хранит как состояние безопасности и более раннюю не принимает. Тяжёлые тела хранятся адресуемо по содержимому, и хэш попадает в журнал только с сертификатом доступности от независимых реплик.

\subsection{Контракт задачи, план и цепочка полномочий}

Контракт задачи фиксирует требуемый результат и правила приёмки: родительскую цель, область записи, входы и зависимости, критерии приёмки, профиль подтверждаемых свойств, класс риска и бюджеты. План отвечает на другой вопрос, что можно сделать: привязки инструментов, гранты и лимиты. Дайджест плана входит в назначение, и среда сверяет его на каждом вызове. Изменение критериев создаёт новую версию контракта с авторизацией того же уровня, что и создание задачи, поэтому продьюсер не может сделать падающий тест зелёным, удалив сам тест. Цепочка полномочий идёт от цели, которую подписывает кор-нода, через полномочие оркестратора к задаче и ревизии, и на каждом уровне права только сужаются. Привилегированный класс риска определяется детерминированно по путям и типам операций (авторизация и доступ, зависимости, сборка и CI, проверки, правила, инструкции агентов и хуки, подпись, публикация, секреты, миграции) и требует двух независимых ревью или одобрения человека из кор-команды. Задачи с пересекающимися областями записи одновременно не исполняются, и область удерживается до интеграции или отмены.

\subsection{Оркестратор по аренде и принятие точного кандидата}

Оркестратор протокола держит роль по аренде с эпохой и случайным fencing-токеном \cite{chubby,kleppmann-fencing}. Каждая его запись несёт эпоху и токен, и валидаторы отвергают запись устаревшей аренды, сохраняя её как расхождение. Принудительный захват роли не передаёт сессию предшественника, и даже планово полученная сессия считается недоверенной: при расхождении прав журнал.

Продьюсер представляет подписанную ревизию $R$ над базой $B$ вместе с головой цепочки квитанций, то есть подписанных средой фактов исполнения его команд. Среда строит неизменяемый \emph{интеграционный кандидат} $I$ над текущей головой $M$ и фиксирует пару $(M,I)$. Множество допустимых ревьюверов замораживается до появления хэша $I$; затем получают случайность только из хэшей, из ``луковицы хэшей'' кор-нод и хэша блока Swarma Chain, вышедшего после фиксации $I$, и ревьювер вычисляется из идентификатора задачи, хэша $I$ и случайности. Оркестратор поэтому не выбирает удобного проверяющего после того, как кандидат известен. Допустимость ревьювера задаёт политика независимости по трём осям: административный домен, семейство модели и тулчейн.

\emph{Аттестацией} называется подписанный результат проверки свойства с исходом ``пройдена'', ``провалена'' или ``не установлена''; таймаут и сбой инструмента дают последний и зелёной проверкой не становятся. Финальное одобрение и каждая аттестация адресованы четвёрке: версия контракта, кандидат $I$, версия политики проверки и дайджест окружения. Проверки инициирует нода, третья по отношению к продьюсеру и ревьюверу, а критические наборы обязательно повторяет верификатор на другой машине. Получив QC на интеграцию, служба действий с репозиторием атомарно заменяет $M$ на $I$, только если текущая голова всё ещё равна $M$. Если голова сдвинулась, кандидат перестраивается, а ревью и проверки повторяются: одинаковый patch-id на другой базе одобрения не переносит.

\subsection{Внешние эффекты и остановка}

Подпись, публикацию и обращение к устройству исполняет раннер с идентификатором и ключом идемпотентности каждой операции и исходом ``выполнена'', ``отклонена'' или ``не установлен''. При неизвестном исходе слепой повтор запрещён: процесс сверяет состояние получателя либо получает решение человека. Авторизация из прежней эпохи или до остановки без новой авторизации не исполняется, и это проверяют демон среды и раннер на границе эффекта. Любая кор-нода может выпустить сертификат экстренной остановки одной подписью, и среда сразу перестаёт открывать области задач, начинать операции и продвигать голову; возобновление требует QC текущего уровня. Один честный участник должен суметь остановить подозрительную работу, и один участник не должен суметь возобновить её вопреки остальным.

\subsection{Федерация и доверие между владельцами}

Федерация объединяет несколько проектов, между которыми контрибьютор сам выбирает получателя ёмкости. Каждый проект публикует подписанный манифест: адрес среды и хэш её демона, уровень доверия, формы клиента, ограничения на ёмкость и провайдеров, инвентарь ограничений среды и политику независимости. При присоединении манифест компилируется в план контрибьютора с показом разницы полномочий. Доверие направленное: контрибьютор доверяет среде и кор-команде выбранного проекта, а проект продолжает считать рабочую ноду недоверенной. Заявка незнакомой ноды несёт PoW-билет допуска, который одновременно служит шаром майнинга Swarma Chain, а ключа подписи допуска на публичной точке присоединения нет.

Люди и агенты разных владельцев переписываются через тот же носитель; белый список серверов и \texttt{m.room.server\_acl} ограничивают федерацию Matrix, а подписи и гранты поверх них проверяются всё равно. Человека другого владельца узнают по его корню, закреплённому один раз при знакомстве. Корень аттестует мастер-ключ перекрёстной подписи человека в Matrix вместе с полным адресом, и устройство отправителя проверяется цепочкой из четырёх звеньев: ключ устройства, ключ самоподписи, мастер-ключ, закреплённый корень. Если звено не сходится, сообщение уровня A1 не получает. Сервер, выдавший адрес, может отправить сообщение от его имени, поэтому гранты называют ключи, удостоверенные корнем, и смена сервера или адреса прав не переносит.

\section{Идентичность, аттестации и репутация}
\label{sec:identity}

Каждый механизм выше называет участника по-своему: UID и поколение на сокете, путь деривации у ключника, адрес на носителе, субъект гранта, издатель записи памяти, ключ ноды в протоколе. Swarma сводит их к одной устойчивой личности, чтобы история участника не зависела от машины, учётки и проекта.

\textbf{Ячейка личности.} Личность задаёт \emph{ячейка личности}, расходуемая запись с тремя полями: текущие ключи, хэш следующих ключей и ссылка на родителя. Идентификатор личности равен хэшу первой записи и при ротациях не меняется. Ячейку меняют три траты: ротация следующими ключами, трата родителем и отзыв в терминальную ячейку. Собственная ротация не меняет ссылку на родителя, а предротация в духе KERI \cite{keri} защищает от поломки текущих ключей, потому что следующие спрятаны за хэшем. Ячейка владельца служит родителем ячеек его агентов; её ключи из мастер-секрета установки не выводятся и хранятся отдельно, а её собственным родителем служит внешний ключ восстановления. Поэтому при компрометации мастер-секрета ячейка владельца переводит ячейки агентов на ключи нового корня, и личности агентов сохраняются.

\textbf{Как всё сходится на одной личности.} Рантайм выдаёт личности и без работающей цепочки: идентификатор самосертифицирующий и равен хэшу первой записи журнала ключевых событий, а журнал ведёт реестр владельца. При публикации в цепочку ячейка получает тот же идентификатор и последнее состояние, и ссылки, накопленные раньше, остаются действительными. На этой ячейке сходятся все представления участника:
\begin{itemize}
\item ключник выводит ключи личности и сообщений агента по стабильному пути, а ключи ячейки выводит из того же мастер-секрета хэш-функцией SHAKE256 от пути и номера ротации;
\item ключ сообщений связан с личностью аттестацией \texttt{attestation/1}, опубликованной в реестре;
\item учётку на носителе выдаёт регистратор по ключу зарегистрированной машины, а мастер-ключ перекрёстной подписи человека аттестует его закреплённый корень;
\item гранты называют ключи, удостоверенные личностью, и адрес в мессенджере субъектом гранта не служит;
\item память называет издателя, субъекта и отправителя записей ячейками и берёт основание представительства из гранта в реестре;
\item в протоколе короткоживущие ключи нод удостоверяет личность участника, а ревизии подписываются ключом ноды в формате SSH-подписи git;
\item аттестации называют выдавшего и субъекта по ячейкам, а имя в Swarma Chain указывает на идентификатор ячейки.
\end{itemize}
Переезд агента на другую машину или в другой рантайм поэтому не рвёт ни истории, ни памяти: меняются ключи и место исполнения, а ячейка остаётся.

\textbf{Доказательство идентичности.} Доказательство связывает личность с внешней сущностью в обе стороны: владелец подписывает заявление ключами ячейки и публикует у внешнего провайдера артефакт с идентификатором, например TXT-запись домена под DNSSEC, файл well-known по HTTPS или ключ аккаунта на хостинге кода. Доказательство удостоверяет контроль над доменом или аккаунтом и честности владельца не удостоверяет. Личности под одной ячейкой владельца, личности, выведенные ключником одной установки, и разные DNS-имена одного владельца относятся к одному административному домену.

\textbf{Аттестации.} После ревью или приёмки рантайм проекта выдаёт аттестацию исхода: личность агента, роль (продьюсер, ревьювер или верификатор), ссылку на проект, хэш кандидата или ревизии, версии контракта и политики, исход ``принято'', ``отклонено'' или ``не установлено'', высоту и время, личность выдавшего и подпись пачки. Выдавать аттестации вправе любой рантайм с личностью. В цепочку попадают только хэши и метаданные; приватный проект публикует хэш-обязательство, и полное свидетельство субъект раскрывает сам. Проверка не требует доверия чужой базе: пачка подписана ключами, текущими у выдавшего на высоте записи, путь Меркла сходится к корню пачки в блоке, ячейка выдавшего жива на этой высоте по корню состояния личностей в заголовке, а вес считается по опубликованной формуле.

\textbf{Репутация.} Репутацию детерминированно вычисляют индексаторы как неподвижную точку графа подписанных утверждений по схеме EigenTrust \cite{eigentrust}. Затравка задаётся классами доказательства личности, оси продьюсера, ревьювера и верификатора считаются раздельно, старые свидетельства весят меньше, а кластер личностей, аттестующих друг друга, получает ровно тот вес, который втекает в него снаружи. Исход ``не установлено'' нейтрален. Спор оформляется подписанным возражением без арбитража. Принимающий рантайм применяет к истории собственную политику: новому агенту можно назначить двух-трёх независимых ревьюверов, агенту с историей одного, но репутация не выдаёт гранта и не создаёт отсутствующей независимости. Пропорциональное вознаграждение при полной устойчивости к множественным личностям недостижимо \cite{sybil,sybilproof}, поэтому каждая личность, заявка и запись стоят работы или платы, а деньги и репутация не связаны.

\section{Swarma Chain: интеграция с платформой}
\label{sec:chain}

Swarma Chain служит нижним слоем общего состояния; её устройство описывает отдельная работа \cite{swarmachain}, а здесь мы называем только точки, где на неё опираются рантайм и протокол. Консенсус цепочки ведёт деньги, имена и состояние личностей, тяжёлый контент лежит вне цепочки в p2p-памяти. У аттестаций и социальных записей узел проверяет конверт (тип, размер, подпись пачки, оплату и статус выдавшего), а смысл детерминированно вычисляют индексаторы.

Консенсус опирается только на хэш-функции. Обязательства строятся как SHA-256~$\|$~SHAKE256 на 64 байта, подписи как пара SLH-DSA-128s в вариантах SHA2 и SHAKE \cite{fips205}, обязательных обеих, весом $2\times7856=15712$ байт. Кривых, пейрингов и решёток в консенсусе нет. Массовые записи подписываются пачками: одна парная подпись покрывает корень дерева Меркла, а доказательство включения носит субъект записи. Заголовок каждого блока несёт корень хэш-дерева состояния личностей и имён, и лёгкий клиент проверяет заголовки от последнего подписанного checkpoint.

Рантайм и протокол берут у цепочки пять вещей. Ячейки личностей делают одну личность общей для всех проектов и снимают предел реестра одного владельца, который может показать разным проверяющим разные истории. Имена указывают на идентификаторы ячеек и регистрируются через обязательство и раскрытие с вступлением в силу через 24 часа. Пачки аттестаций выносят историю работы за границы журналов проектов. Хэш блока, вышедшего после фиксации кандидата, входит в случайность выбора ревьювера, а PoW-билеты допуска публичных нод служат шарами майнинга. Наконец, головы журналов протокола и реестров владельцев закрепляются в цепочке, и запись, наблюдённая до поломки класса ключей, сохраняет доказательную силу на дату наблюдения.

Вне консенсуса корни доверия стоят на кривых: Ed25519 в журнале протокола, квитанциях и ключах рантайма, P-256 в защищённых чипах телефонов, SSH, TLS и подписи git. План миграции на случай поломки кривых состоит из четырёх шагов: история фиксируется с датой наблюдения, долгоживущие корни переходят на хэш-подписи, а ключи на кривых остаются короткоживущими листьями под ними, сеансовый слой следует за вендорами с гибридным обменом ML-KEM, а поломанный класс ключей отвергается политикой, и личности ротируются через предротацию.

\section{Анализ безопасности по слоям}
\label{sec:analysis}

Гарантии ниже названы вместе с противником, механизмом, допущением и пределом (табл.~\ref{tab:guarantees}). Это аргументы уровня дизайна; подробный разбор каждой гарантии рантайма с доказательными шагами дан в работе-спутнике \cite{swarmatrust}, а протокольные инварианты разобраны в книге \cite{swarmabook}. Общее условие для всех строк: гарантия держится, только если исполнитель эффекта вызывает проверку и физически удерживает область эффекта.

\begin{table}[t]
\centering
\caption{Гарантии по слоям. П1--П8 обозначают противников из раздела~\ref{sec:model}.}
\label{tab:guarantees}
\footnotesize
\begin{tabular}{L{3.3cm}L{1.15cm}L{3.9cm}L{5.95cm}}
\toprule
Гарантия & Против & Механизм & Допущение и предел \\
\midrule
G1. Агент не получает закрытого ключа своей личности & П1 & ключник, UID от ядра, активное поколение & ядро, root и мастер-секрет целы; захваченный агент получает подписи своей личности, и подпись правоты не доказывает \\
\addlinespace
G2. Делегирование не расширяет полномочий & П1, П2 & ослабление, ключ агента не издаёт звеньев, вызывающий в цепочке & вредное действие в пределах гранта остаётся остаточным риском \\
\addlinespace
G3. Отзыв действует не позже $\max(S+A,F)$, неизвестность закрывает допуск & П1, П5 & накопительные снимки, записи свежести, три исхода проверки & ключ реестра цел; токен, уже выданный процессу, пригоден до аннулирования у провайдера \\
\addlinespace
G4. Откат и раздвоение журнала реестра обнаруживаются & П5 & нумерованные контрольные точки, свидетель, точка вне узла & реестр и свидетель на одном узле образуют один домен отказа; разные истории перед разными проверяющими снимает только цепочка \\
\addlinespace
G5. Доставка, членство в комнате и подпись коллеги полномочий не дают & П2, П3 & ворота доверия, звонок и почтовый ящик, лестница A0--A3 & decision guard инъекцию надёжно не распознаёт \cite{agentdojo}; при политике A0 компрометация сервера даёт команды в зоне человека \\
\addlinespace
G6. Учётки на носителе ограничены меткой машины & П1, П4 & регистратор, пропуск 120~с, токен без продления & захваченный агент получает пропуск на своё разрешённое имя \\
\addlinespace
G7. Исполнение только на станции с квалифицированными гарантиями & П1 & контракт изоляции, квитанция границы, отказ на \texttt{absent} и \texttt{PENDING-VERIFY} & самоотчёт машины не защищает от захваченного root и скопированного ключа машины; CONNECT оставляет канал выноса \\
\addlinespace
G8. Процесс не принимает изменение без независимого ревью точной ревизии & П1, П2 & правила ядра доски, свежесть свидетельства, предикат над накопленной дельтой & ревьювер может ошибиться; общая модель коррелирует ошибки \cite{selfpref} \\
\addlinespace
G9. Принятие относится к точному кандидату, эффект только по актуальной записи & П2, П6, П7 & QC, аренда с fencing, четвёрка адреса, CAS, выбор ревьювера после фиксации $I$ & среда доверена в целостности; на L1 единственность истории держит секвенсер \\
\addlinespace
G10. Память и репутация полномочий не создают & П1, П2 & нет операций выдачи, недоступный реестр не грант, политика у получателя & история удачных работ говорит только о прошлом; из истории git ничего не стирается \\
\bottomrule
\end{tabular}
\end{table}

\textbf{Захваченный оркестратор.} В рантайме захваченный оркестратор пользуется ключом вызывающего из своего дома, но силу этого ключа ограничивает его грант, и отзыв гранта отрезает созданных им агентов. Грант общего языка он соседу не выдаст: издатель звена с аттестацией агента отклоняется, даже если подписи верны, а суженный грант брокера v0 для созданного им оркестратора отрезается вместе с его собственным при отзыве. В протоколе у оркестратора нет прямого shell среды; он действует записями, валидаторы проверяют область, лимиты, класс риска и независимость, ревьювера вычисляет правило выбора, а любая кор-нода может остановить среду одной подписью. Остаточный риск остаётся там, где правильно исполненную вредную задачу приняли до обнаружения; его ограничивают область полномочия, бюджет, отложенное закрепление начислений и остановка.

\textbf{Захваченная машина (П4).} Root на машине с ключником получает мастер-секрет установки или возможность им пользоваться, а с ним ключи всех агентов установки и их гранты до отзыва. Ключ оператора на телефоне, ключ реестра на узле доверия и офлайновый издатель регистраций остаются вне машины: противник не подписывает действий оператора, не подделывает контрольных точек и не регистрирует новых машин, а его учётки на носителе ограничены именами этой машины. Восстановление идёт через отзыв машины и новый корень, а ячейка владельца переводит ячейки агентов на новые ключи.

\textbf{Среда протокола и криптографический противник.} Скомпрометированная среда может показать один diff и интегрировать другой или сообщить вымышленный результат проверок, и протокол этого не предотвращает: квитанция среды доказывает авторство отчёта, а расхождение обнаруживает только верификатор на другой машине. Против П8 гарантии вне Swarma Chain держатся лишь в пределах плана миграции из раздела~\ref{sec:chain}, а предротация ячеек не даёт противнику перехватить личность через поломанные текущие ключи.

\section{Оценка}
\label{sec:eval}

Оценка сообщает измеренные факты об эталонных компонентах платформы: согласие независимых реализаций на конечных корпусах, привилегированные шаги хоста станции на машине Linux, исчерпывающий перебор ограниченной модели интеграции, стоимость эталонной проверки записей и статистику процесса ревью (табл.~\ref{tab:evidence}). Задержка отзыва на сквозном пути станции, производительность системы целиком, состязательная проверка живой установки и польза роя на внешних задачах в эту оценку не входят; методы последних двух описаны в книге \cite{swarmabook}.

\textbf{Доступность артефактов.} Библиотека проверки подписанных записей с корпусом опубликована в организации \href{https://github.com/relux-works}{github.com/relux-works}; корпус проверен на ревизии f76a400, замеры сделаны на её предке 9b42338. Пакет спецификации памяти лежит в каталоге \texttt{spec/} репозитория \href{https://github.com/relux-works/swarma-memory}{relux-works/swarma-memory} (ревизия 763e024), запись о привилегированных запусках находится в \texttt{docs/root-qualification-evidence.md} репозитория \href{https://github.com/relux-works/swarma-user-manager}{relux-works/swarma-user-manager}. Модель интеграции приложена к исходникам работы-спутника \cite{swarmatrust}. Обвязка замеров, отчёт о замерах и записи ревью не опубликованы.

\begin{table}[t]
\centering
\caption{Сводка свидетельств. Каждая строка относится к названной ревизии и названным условиям.}
\label{tab:evidence}
\footnotesize
\begin{tabular}{L{3.2cm}L{4.4cm}L{4.1cm}L{3.0cm}}
\toprule
Утверждение & Метод & Результат & Чего не показывает \\
\midrule
Генератор и две независимые проверки подписанных записей согласуются & корпус соответствия: генератор на Go, независимая проверка на Python, библиотека на Go; у отказа точный код и положительный контроль & 2632 случая (2499 отрицательных, 133 положительных); согласие трёх сообщает репозиторий, тест библиотеки повторён нами & однозначности и полноты спецификации, работы проверки внутри служб \\
\addlinespace
Два вычислителя ядра памяти согласуются & один вычислитель по сводному контракту без доступа к пакету, второй по пакету & ожидаемые ответы на всех 66 смысловых векторах корпуса 0.2.7 у обоих; на 77 векторах 0.3.0 у вычислителя на Go & защищённой службы, пользы памяти \\
\addlinespace
Привилегированные шаги хоста исполнимы на Linux & два запуска под root на одной машине, привязанные к сборкам & примитивы учёток: применение, проверка, откат; квалификационный набор: 11 из 11 проверок при применении и при проверке & самоквалификации по 42 гарантиям, перезагрузок, других машин \\
\addlinespace
Четыре правила доставки необходимы для свойства интеграции & исчерпывающий перебор ограниченной модели с отрицательным контролем на каждое правило & 1090 состояний без нарушения; без любого правила найден контрпример & неограниченных параметров, эквивалентности реализации \\
\addlinespace
Стоимость эталонной проверки одной записи & микрозамеры библиотеки на входах корпуса в одном потоке Go & медианы от 2,1 до 18,0~мс на вызов (табл.~\ref{tab:bench}) & задержки служб и сети, оптимизированной реализации \\
\addlinespace
Ревьювер возвращает заметную долю первых ревизий & разбор записей ревью на доске задач одного проекта & приняты 17 из 35 разобранных первых раундов; медиана 2 раунда на задачу & пропущенных дефектов, причинного эффекта \\
\bottomrule
\end{tabular}
\end{table}

\subsection{Корпуса соответствия}

В корпусе подписанных записей каждый случай отказа несёт точный ожидаемый код и ссылку на положительный контроль, поэтому реализация, которая отказывает всегда, корпус не проходит. Генератор на Go пишет случаи и ожидания, проверка на Python сама канонизирует байты и проверяет подписи Ed25519 и P-256, вердиктов Go не импортируя. В репозитории лежат 944 именованных случая, а 1688 механических развёрток материализуются по манифесту с хэшами SHA-256. Тест библиотеки генерирует полный корпус, сверяет манифест и требует точного кода и плана для всех 2632 случаев. Полный прогон проверки на Python сообщает репозиторий; тест библиотеки мы повторили на ревизии f76a400 (Go~1.26.0, darwin/arm64): 2632 подтеста прошли, код выхода 0. Процедура корпуса задаёт и отрицательный контроль самих инструментов: изменённый байт подписи в положительном случае обязан дать расхождение у генератора и код \texttt{envelope\_signature\_invalid} у проверки на Python.

Для памяти первый вычислитель написан по сводному контракту ядра без доступа к пакету и прогнан на векторах пакета без их правки, второй написан только по пакету, и каждое расхождение разобрано как ошибка пакета, реализации или неоднозначность контракта. На корпусе 0.2.7 оба дают ожидаемые ответы на всех 66 смысловых векторах; на корпусе 0.3.0 вычислитель на Go проходит все 77 смысловых векторов. Восемь векторов защищённого развёртывания этими вычислителями не исполняются.

\subsection{Привилегированные шаги хоста станции}

На одной машине Linux выполнены два привилегированных запуска, каждый на своей сборке. Первый прошёл режимы применения, проверки и отката для одноразовых примитивов учёток: команды учёток и групп, защита дома, отказ ACL и capabilities, остановка через pidfd. Второй установил набор квалификации: при применении квалификация в режиме загрузки прошла 11 из 11 проверок, отдельная проверка тоже прошла 11 из 11, а допуск станций остался не настроен. Опубликованная запись о результатах указывает, что второй запуск не устанавливает ни кампании перезагрузок, ни отката установки. Набор из 11 проверок не совпадает с самоквалификацией по 42 гарантиям из раздела~\ref{sec:runtime} и подписанной аттестации машины не образует.

\subsection{Ограниченная модель интеграции}

Модель воспроизводится командой \texttt{make -C paper model} из исходников работы-спутника. Она содержит две задачи, до двух версий контракта, до двух эпох аренды и до двух доставок записи и на границе доставки применяет четыре правила: голова равна базе записи (\texttt{cas}), запись несёт текущую эпоху (\texttt{epoch}), аттестация ссылается на текущий контракт (\texttt{contract}) и адресована точному кандидату (\texttt{bind\_I}). При всех правилах перебор в ширину достигает 1090 состояний без нарушения, а отключение любого одного правила даёт контрпример (табл.~\ref{tab:model}); числа отмечают остановку поиска на первом нарушении и безопасность вариантов не сравнивают. Модель считает среду и валидаторов честными и не описывает содержимого кода, полной аренды, остановки и внешних эффектов.

\begin{table}[t]
\centering
\caption{Ограниченная модель интеграции: напечатанный результат каждого запуска.}
\label{tab:model}
\footnotesize
\begin{tabular}{L{2.6cm}rL{7.8cm}r}
\toprule
Запуск & Состояний & Напечатанное нарушение & Действий в трассе \\
\midrule
Все правила & 1090 & \texttt{None} & нет \\
Без \texttt{cas} & 84 & \texttt{record applied twice} & 6 \\
Без \texttt{epoch} & 82 & \texttt{record from a stale lease epoch advanced main} & 6 \\
Без \texttt{contract} & 71 & \texttt{attestation references a superseded contract version} & 6 \\
Без \texttt{bind\_I} & 490 & \texttt{attestation made on a different base than the integrated candidate} & 11 \\
\bottomrule
\end{tabular}
\end{table}

\subsection{Микрозамеры эталонной проверки записей}

Мы измерили шесть публичных вызовов библиотеки проверки на ревизии 9b42338, на неизменённых входах корпуса; реализация эталонная и не оптимизировалась. Замер шёл в гостевой системе KVM на процессоре AMD Ryzen 7 9700X (гостю видны 8 логических CPU, без закрепления), ядро Linux 7.0.0-38-generic, Go 1.26.10 linux/amd64, \texttt{GOMAXPROCS=1}, пять последовательных повторов стандартного бенчмарка Go с \texttt{-benchmem}. Машина была общей: одноминутная средняя нагрузка составляла 1,73 до замера и 1,95 после. Внутри таймера библиотека разбирает и проверяет сырые подписанные байты, а загрузка фикстур, ввод-вывод и подпись в замер не входят.

\begin{table}[t]
\centering
\caption{Эталонная проверка записей на вызов. Время дано медианой и размахом [минимум; максимум] пяти повторов; размах не является доверительным интервалом. Память ($10^6$ байт) и число выделений даны медианой.}
\label{tab:bench}
\footnotesize
\begin{tabular}{L{7.6cm}rrr}
\toprule
Вызов и вход & Время, мс & Память & Выделений \\
\midrule
Сообщение \texttt{message/1}, 835 байт & 6,04 [5,70--7,22] & 6,30 & 92\,619 \\
Грант, одно звено, 832 байта & 9,10 [7,76--10,53] & 9,00 & 132\,307 \\
Грант, три звена, 2628 байт & 17,99 [15,53--23,75] & 18,91 & 278\,351 \\
Аттестация ключа сообщений, 840 байт & 3,78 [3,71--4,51] & 4,49 & 66\,016 \\
Снимок отзывов на 64 записи, 7110 байт & 2,11 [1,67--2,72] & 1,94 & 28\,497 \\
Журнал из 64 записей и две контрольные точки, 30\,718 байт & 9,21 [7,54--11,23] & 7,53 & 112\,915 \\
\bottomrule
\end{tabular}
\end{table}

Разница одного и трёх звеньев не оценивает влияния глубины: входы различаются содержимым, и в трёхзвенной цепочке есть подпись P-256; худший поиск среди действующих отзывов не измерен. При медиане 6,04~мс на сообщение один поток проверяет около 165 сообщений в секунду, и для ворот доверия это стоимость эталонной, неоптимизированной реализации.

\subsection{Статистика ревью}

Разработка проекта Swarma Curator идёт на доске задач с продьюсером и независимым ревьювером из раздела~\ref{sec:workflow}, и ревьювер выносит вердикт с находками по классам серьёзности. Мы разобрали файлы исходов ревью на доске задач этого проекта по состоянию на 24 октября 2026 года. Вердиктом считалось последнее поле \texttt{verdict} в файле, а серьёзности брались из JSON находок. Нашлось 140 задач с записями ревью и 318 записей. Вердикт разобран в 104 записях из 318; остальные 214 написаны в ранних форматах, поэтому выборка смещена к поздним записям. Из 104 разобранных вердиктов 55 \emph{принято} и 49 \emph{нужны изменения}; из 35 разобранных первых раундов принято 17. Медиана числа раундов на задачу равна 2, максимум 13. Находок высшего класса серьёзности P0 нет, классов P1 и P2 соответственно 73 и 49. Второй разборщик текстовых вердиктов даёт 58 и 51 при 209 неразобранных записях. Числа задач, записей и находок служат нижней границей для этой доски; медиана, доля принятых первых раундов и соотношение вердиктов нижней границей не являются, и пропущенные записи могут сдвинуть их в любую сторону. Числа показывают, что ревьювер регулярно возвращает работу, и ничего не говорят о дефектах, которые он пропустил.

\section{Смежные работы}
\label{sec:related}

\textbf{Агентские системы.} ReAct \cite{react} связал рассуждение модели с вызовом инструментов, SWE-agent \cite{sweagent} и SWE-bench \cite{swebench} перенесли это на задачи разработки, AutoGen, MetaGPT и ChatDev \cite{autogen,metagpt,chatdev} организуют несколько ролей и обмен между ними. Параллельные агенты, написавшие компилятор \cite{compiler}, опирались на тесты как оракул, а в Prove2Me \cite{prove2me} работу агентов независимых владельцев принимает ядро доказательств. Эти работы изучают способности и координацию; Swarma добавляет ключи, полномочия, отзыв и процедуру принятия кода без механического оракула.

\textbf{Системная защита агентов от инъекций.} Инъекция против агентов с инструментами описана в \cite{greshake} и измерена в AgentDojo \cite{agentdojo}. Шаблон двух моделей \cite{duallm} не показывает недоверенный текст модели с правом вызывать инструменты, CaMeL \cite{camel} извлекает поток управления из доверенного запроса и проверяет политику на capabilities данных, FIDES \cite{fides} отслеживает метки конфиденциальности и целостности, Progent \cite{progent} задаёт агенту политику привилегий, IsolateGPT \cite{isolategpt} исполняет приложения на базе модели изолированно, а \cite{designpatterns} сводит такие решения в шаблоны. С ними Swarma разделяет два принципа: допуск решает детерминированный код, а оценка моделью только сужает. Эти работы ограничивают потоки данных и вызовы внутри одного агентского приложения; Swarma ставит проверку у исполнителя эффекта за границей процесса, между машинами и владельцами, и их механизмы применимы внутри станции как дополнительное сужение.

\textbf{Полномочия и прозрачность.} Гранты продолжают линию capability-систем \cite{dennis1966,miller2006}; Macaroons \cite{macaroons} добавляют ограничения без обращения к издателю, Biscuit \cite{biscuit} даёт офлайн-ослабление с языком правил, SPKI/SDSI \cite{spki} и UCAN \cite{ucan} строят цепочки делегирования на открытых ключах. Certificate Transparency \cite{ct,ct2}, CONIKS \cite{coniks} и Sigstore \cite{sigstore} делают раздвоение журналов обнаружимым, а реестр владельца решает ту же задачу для одного владельца; журнал ключевых событий и предротация ячейки личности следуют KERI \cite{keri}.

\textbf{Изоляция, мессенджеры, консенсус и репутация.} gVisor \cite{gvisor}, Firecracker \cite{firecracker} и Kata Containers \cite{kata} изолируют недоверенный код ядром в пространстве пользователя или лёгкими виртуальными машинами, а станция входит в тот же ряд через контракт изоляции. Перекрёстная подпись Matrix \cite{matrix} и MLS \cite{mls} удостоверяют устройства; Swarma использует перекрёстную подпись как уровень A1 и добавляет подпись точных байтов и гранты. Журнал протокола следует BFT-протоколам и кворумным системам \cite{pbft,hotstuff,byzquorum}, аренда с fencing-токеном следует \cite{chubby,kleppmann-fencing}, а расчёт репутации и его пределы опираются на \cite{eigentrust,sybil,sybilproof}. Автоматизация ревью кода изучалась в \cite{tufano2021,tufano2022,codereviewer,frommgen2024}, а смещение моделей-судей, в том числе к собственному тексту \cite{llmjudge,selfpref}, объясняет требование другого издателя модели для ревьювера.

\section{Заключение}

Swarma переносит решение о допуске из суждения модели в детерминированные проверки у исполнителя эффекта и делает это на всех трёх слоях, а процесс работы оставляет командам как подключаемый модуль над неизменными правилами ядра. Цена этого переноса видна в допущениях: доверие сводится к ядру, root и доверенной базе хоста, к целому мастер-секрету ключника, к проверенному источнику времени, к узлу доверия под администрированием владельца, к человеку с аппаратным ключом и, в протоколе, к целостности среды проекта.

Измеренные свидетельства относятся к эталонным компонентам: корпусу подписанных записей, вычислителям памяти, привилегированным шагам хоста на Linux, ограниченной модели интеграции, стоимости проверки записей и процессу ревью. Свойства сквозного пути станции, то есть задержка отзыва через реестр, брокер, диспетчер и хост сессий и поведение программных нарушителей на настоящих границах, а также польза роя на одинаковых внешних задачах требуют отдельных измерений. Их методы, включая калибровку ревьюверов посевом дефектов, описаны в книге \cite{swarmabook}, а аренда, остановка и внешние эффекты протокола допускают механически проверяемую спецификацию \cite{tla}.

\raggedcolumns
\begin{multicols}{2}
\footnotesize\raggedright
\bibliographystyle{abbrv}
\setlength{\itemsep}{0pt}
\bibliography{references}
\end{multicols}

\end{document}
